\section{Operational lessons}
\label{sec:appendix-operational}

Running TML-bench reliably surfaced several operational lessons that may be useful to teams building similar agent evaluation systems.

\subsection{A control plane beats ad-hoc scripts}
Long-running suites benefit from a simple ``control plane'': durable run IDs, structured logs, and machine-readable events. This makes suites resumable, debuggable, and auditable. In contrast, ad-hoc shell scripts tend to fail silently, make partial failures hard to diagnose, and are difficult to parallelize safely.

\subsection{Circuit breakers save time and money}
When a provider or model enters a failure streak (timeouts, repeated invalid submissions, intermittent API errors), a circuit breaker can pause or skip that lane. This prevents burning compute and wall-clock on runs that are likely to fail again and allows the suite to make progress elsewhere.

\subsection{Incremental persistence prevents ``lost suites''}
Suites that take hours or days should write results continuously. Persisting every run outcome (including failures) as soon as it completes makes the overall process robust to interruptions: machine restarts, parent process crashes, or transient provider outages. It also supports incremental analysis rather than ``all-or-nothing'' end-of-suite reporting.

\subsection{Per-task resource caps improve stability}
Some tasks are more resource-sensitive than others (for example, heavy preprocessing, large feature matrices, or memory pressure under parallelism). Per-task caps (concurrency limits, memory limits, or stricter runtime limits) can prevent cascading failures and make the benchmark safer to run repeatedly.

\subsection{Post-run diagnostics are part of reliability}
A benchmark should treat ``why did this run fail?'' as a first-class output. Capturing a compact post-run diagnostic artifact (status, timeout vs.\ validation vs.\ runtime error, and a short trace or log pointer) turns failures into actionable debugging items and improves reproducibility of observed failure modes.
